\documentclass[12pt]{article}
\usepackage[utf8]{inputenc}
\usepackage[T1]{fontenc}
\usepackage{hyperref}
\title{Muzi Language Resources for K-12 Curriculum Development}
\author{Compiled by Claude Sonnet 5.5 (claude-sonnet-5-5) for Andr\'as Kornai}
\date{}
\begin{document}
\maketitle

\section{Introduction}
Muzi (ISO 639-3 ymz; Glottocode muzi1235) is a Loloish (Sino-Tibetan) language of Yunnan, China, associated with the Phula people. Wikipedia gives about 10,000 speakers (2008) and lists communities in Gejiu and Jinping counties. It was split from the former code ymj in 2008 and should not be confused with the other Muji languages. Search found only encyclopedic and catalogue entries.

\section{Primary Educational Resources}
No primers, textbooks or school programmes were found. The language is not known to be taught in schools.

\section{Linguistic and Reference Materials}
\begin{itemize}
\item Wikipedia, Muzi language: \url{https://en.wikipedia.org/wiki/Muzi_language}
\item Wikipedia, Southeastern Loloish languages (comparative tables): \url{https://en.wikipedia.org/wiki/Southeastern_Loloish_languages}
\item Wikipedia, Loloish languages: \url{https://en.wikipedia.org/wiki/Loloish_languages}
\item Glottolog, languoid page for muzi1235: \url{https://glottolog.org/resource/languoid/id/muzi1235.bigmap.html}
\end{itemize}

\section{Religious and Community Resources}
None found.

\section{Digital and Audio Resources}
None found.

\section{Curriculum Development Context}
No usable curriculum material exists in what was found. Documentation (word lists, recordings, an orthography) would have to come first, and the Glottolog and Wikipedia references can serve as starting points for locating the underlying academic sources, such as the Pelkey (2011) work cited there.

\end{document}